\documentclass[aspectratio=169]{beamer}
\usepackage[utf8]{inputenc}
\usepackage{xeCJK}
\usepackage{graphicx}
\usepackage{booktabs}
\usepackage{tikz}
\usepackage{pgfplots}
\usepackage{xcolor}

% 設定中文字體
\setCJKmainfont{Noto Sans CJK TC}

% 主題設定
\usetheme{Madrid}
\usecolortheme{seahorse}

% 標題頁資訊
\title{CBME與EPA在心臟內科的應用}
\subtitle{以能力為導向的醫學教育實務}
\author{內科部/心臟內科/謝慕揚醫師}
\date{\today}
\institute{醫學中心- 新竹台大分院}

% 自訂顏色
\definecolor{cardioBlue}{RGB}{0,82,155}
\definecolor{lightGray}{RGB}{240,240,240}

\begin{document}

% 標題頁
\begin{frame}
\titlepage
\end{frame}

% 大綱
% \begin{frame}{今日議程}
% \tableofcontents
% \end{frame}

\section{什麼是EPA？}

\begin{frame}{EPA基本概念}
\begin{block}{Entrustable Professional Activities (EPA)}
\textbf{可信賴專業活動} - 醫師在臨床實務中能夠獨立勝任的專業工作
\end{block}

\vspace{0.5cm}

\begin{columns}[c]
\column{0.5\textwidth}
\textbf{傳統教育模式：}
\begin{itemize}
\item 以時間為基礎
\item 知識導向
\item 標準化評估
\end{itemize}

\column{0.5\textwidth}
\textbf{EPA導向模式：}
\begin{itemize}
\item 以能力為基礎
\item 實務導向
\item 個別化進展
\end{itemize}
\end{columns}
\end{frame}

\begin{frame}{EPA的核心特徵}
\begin{tikzpicture}[remember picture, overlay]
\node[anchor=center] at (current page.center) {
\begin{tikzpicture}
% 中心圓圈
\fill[cardioBlue] (0,0) circle (1.5cm);
\node[white, font=\Large\bfseries] at (0,0) {EPA};

% 外圍特徵
\foreach \angle/\text in {0/可觀察, 60/可評估, 120/可信賴, 180/專業性, 240/整合性, 300/實務性} {
  \node[circle, fill=lightGray, minimum size=2cm, font=\small] 
       at (\angle:3cm) {\text};
  \draw[thick, cardioBlue] (0,0) -- (\angle:1.8cm);
}
\end{tikzpicture}
};
\end{tikzpicture}
\end{frame}

\section{為什麼需要EPA？}

\begin{frame}{醫學教育面臨的挑戰}
\begin{alertblock}{現況問題}
\begin{itemize}
\item 理論與實務脫節
\item 評估方式單一化
\item 學習進度標準化
\item 能力認證不明確
\end{itemize}
\end{alertblock}

\vspace{0.5cm}

\begin{exampleblock}{EPA的解決方案}
\begin{itemize}
\item 整合性能力評估
\item 個別化學習路徑
\item 真實情境評估
\item 漸進式責任委託
\end{itemize}
\end{exampleblock}
\end{frame}

\begin{frame}{國際趨勢與發展}
\begin{center}
\begin{tikzpicture}
\begin{axis}[
    ybar,
    xlabel={年份},
    ylabel={採用EPA的醫學院數量},
    ymin=0,
    ymax=500,
    symbolic x coords={2010,2015,2020,2025},
    xtick=data,
    nodes near coords,
    nodes near coords align={vertical},
    width=10cm,
    height=6cm,
    bar width=30pt,
]
\addplot coordinates {(2010,20) (2015,150) (2020,350) (2025,450)};
\end{axis}
\end{tikzpicture}
\end{center}

\footnotesize{*資料來源：國際醫學教育聯盟報告}
\end{frame}

\section{心臟內科EPA框架}

\begin{frame}{心臟內科10大EPAs概覽}
\scriptsize
\begin{table}
\centering
\begin{tabular}{p{0.1\textwidth}p{0.4\textwidth}p{0.35\textwidth}}
\hline
\textbf{EPA} & \textbf{專業活動} & \textbf{核心能力} \\
\hline
EPA 1 & 病史詢問與風險評估 & 症狀評估、危險因子識別 \\
\hline
EPA 2 & 心血管理學檢查 & 聽診、血管檢查、水腫評估 \\
\hline
EPA 3 & 心電圖判讀分析 & 心律、ST-T變化、傳導異常 \\
\hline
EPA 4 & 心臟超音波操作 & 影像獲取、功能評估 \\
\hline
EPA 5 & 急性冠心症處理 & 快速診斷、風險分層、治療 \\
\hline
EPA 6 & 心律不整診療 & 心律識別、治療選擇 \\
\hline
EPA 7 & 心衰竭管理 & 診斷評估、藥物調整 \\
\hline
EPA 8 & 高血壓管理 & 診斷、治療、生活介入 \\
\hline
EPA 9 & 心導管判讀 & 影像判讀、治療建議 \\
\hline
EPA 10 & 照護協調溝通 & 團隊合作、病患衛教 \\
\hline
\end{tabular}
\end{table}
\end{frame}

\begin{frame}{里程碑發展架構}
\begin{center}
\begin{tikzpicture}
% 定義顏色
\definecolor{level1}{RGB}{255,182,193}
\definecolor{level2}{RGB}{255,218,185}
\definecolor{level3}{RGB}{255,255,224}
\definecolor{level4}{RGB}{144,238,144}
\definecolor{level5}{RGB}{173,216,230}

% 繪製發展階梯
\fill[level1] (0,0) rectangle (2,1);
\node at (1,0.5) {\scriptsize Level 1};
\node at (1,-0.5) {\scriptsize Novice};

\fill[level2] (2,0) rectangle (4,1);
\node at (3,0.5) {\scriptsize Level 2};
\node at (3,-0.5) {\scriptsize Advanced};
\node at (3,-0.7) {\scriptsize Beginner};

\fill[level3] (4,0) rectangle (6,1);
\node at (5,0.5) {\scriptsize Level 3};
\node at (5,-0.5) {\scriptsize Competent};

\fill[level4] (6,0) rectangle (8,1);
\node at (7,0.5) {\scriptsize Level 4};
\node at (7,-0.5) {\scriptsize Proficient};

\fill[level5] (8,0) rectangle (10,1);
\node at (9,0.5) {\scriptsize Level 5};
\node at (9,-0.5) {\scriptsize Expert};

% 邊框
\draw[thick] (0,0) rectangle (2,1);
\draw[thick] (2,0) rectangle (4,1);
\draw[thick] (4,0) rectangle (6,1);
\draw[thick] (6,0) rectangle (8,1);
\draw[thick] (8,0) rectangle (10,1);

% 箭頭
\draw[->, thick] (0,-1) -- (10,-1);
\node at (5,-1.5) {能力發展進程};

% 標示重點
\node[above] at (2,1.2) {\scriptsize 需督導};
\node[above] at (6,1.2) {\scriptsize 獨立作業};
\node[above] at (9,1.2) {\scriptsize 指導他人};
\end{tikzpicture}
\end{center}

\vspace{0.5cm}
\begin{itemize}
\item \textbf{Level 1-2}：在督導下學習與執行
\item \textbf{Level 3}：能獨立勝任標準作業
\item \textbf{Level 4-5}：能處理複雜情況並指導他人
\end{itemize}
\end{frame}

\section{實施方式}

\begin{frame}{月度晨會Mini-OSCE整合}
\begin{block}{時間規劃 (60分鐘)}
\begin{itemize}
\item 準備說明：5分鐘
\item 三站評估：45分鐘 (3輪 × 15分鐘)
\item 總結回饋：10分鐘
\end{itemize}
\end{block}

\begin{columns}[c]
\column{0.5\textwidth}
\textbf{月度主題安排：}
\begin{itemize}
\item 第1個月：基礎評估技能
\item 第2個月：診斷與急症處理
\item 第3個月：慢性疾病管理
\end{itemize}

\column{0.5\textwidth}
\textbf{評估方式：}
\begin{itemize}
\item 6位學生，2人一組
\item 3個評估站輪轉
\item 即時口頭回饋
\item 書面評估報告
\end{itemize}
\end{columns}
\end{frame}

\begin{frame}{評估工具與方法}
\begin{table}
\centering
\scriptsize
\begin{tabular}{lll}
\hline
\textbf{EPA類型} & \textbf{評估工具} & \textbf{應用範例} \\
\hline
技能操作類 & DOPS & 心臟超音波操作 \\
\hline
臨床推理類 & Mini-CEX & 急性冠心症診斷 \\
\hline
複雜決策類 & CbD & 心導管結果判讀 \\
\hline
溝通協調類 & MSF & 病患照護協調 \\
\hline
緊急處理類 & 模擬評估 & 心律不整急救 \\
\hline
\end{tabular}
\end{table}

\vspace{0.5cm}
\begin{alertblock}{關鍵成功因素}
\begin{itemize}
\item 評估者訓練與標準化
\item 即時且具體的回饋
\item 學習歷程持續追蹤
\end{itemize}
\end{alertblock}
\end{frame}

\section{EPA的使用者}

\begin{frame}{EPA的主要使用者群體}
\begin{center}
\begin{tikzpicture}
% 中心EPA
\node[circle, fill=cardioBlue, text=white, font=\Large\bfseries, minimum size=2cm] at (0,0) {EPA};

% 四個主要使用者群體
\node[rectangle, fill=red!20, text=black, minimum width=3cm, minimum height=1.5cm] at (-3,2) {
\begin{tabular}{c}
\textbf{學習者} \\
醫學生 \\
住院醫師 \\
專科醫師
\end{tabular}
};

\node[rectangle, fill=green!20, text=black, minimum width=3cm, minimum height=1.5cm] at (3,2) {
\begin{tabular}{c}
\textbf{教育者} \\
主治醫師 \\
住院醫師 \\
臨床教師
\end{tabular}
};

\node[rectangle, fill=blue!20, text=black, minimum width=3cm, minimum height=1.5cm] at (-3,-2) {
\begin{tabular}{c}
\textbf{評估者} \\
CCC委員 \\
考試委員 \\
品質管控
\end{tabular}
};

\node[rectangle, fill=orange!20, text=black, minimum width=3cm, minimum height=1.5cm] at (3,-2) {
\begin{tabular}{c}
\textbf{管理者} \\
教育主管 \\
醫院管理 \\
學會組織
\end{tabular}
};

% 連接線
\foreach \angle in {45, 135, 225, 315} {
  \draw[thick, cardioBlue] (0,0) -- (\angle:2.2cm);
}
\end{tikzpicture}
\end{center}
\end{frame}

\begin{frame}{各使用者的應用重點}
\begin{columns}[c]
\column{0.5\textwidth}
\textbf{學習者觀點：}
\begin{itemize}
\item 明確學習目標
\item 自我評估工具
\item 能力發展路徑
\item 個人學習計畫
\end{itemize}

\textbf{教育者觀點：}
\begin{itemize}
\item 教學目標設定
\item 評估標準依據
\item 回饋框架
\item 課程設計指引
\end{itemize}

\column{0.5\textwidth}
\textbf{評估者觀點：}
\begin{itemize}
\item 能力認證標準
\item 評估工具設計
\item 品質保證機制
\item 進展追蹤系統
\end{itemize}

\textbf{管理者觀點：}
\begin{itemize}
\item 課程品質管控
\item 資源配置規劃
\item 政策制定參考
\item 認證標準建立
\end{itemize}
\end{columns}
\end{frame}

\section{實施效益與展望}

\begin{frame}{EPA實施的預期效益}
\begin{block}{教育品質提升}
\begin{itemize}
\item 學習目標更明確具體
\item 評估方式更多元實用
\item 個別化學習支持
\item 能力認證更客觀
\end{itemize}
\end{block}

\begin{block}{臨床照護改善}
\begin{itemize}
\item 醫師能力標準化
\item 病人安全保障
\item 醫療品質提升
\item 團隊協作強化
\end{itemize}
\end{block}
\end{frame}

\begin{frame}{實施時程規劃}
\begin{center}
\begin{tikzpicture}
% 時間軸
\draw[thick] (0,0) -- (10,0);
\foreach \x in {0,2.5,5,7.5,10} {
  \draw (\x,-0.1) -- (\x,0.1);
}

% 時程標示
\node[below] at (0,-0.2) {第1季};
\node[below] at (2.5,-0.2) {第2季};
\node[below] at (5,-0.2) {第3季};
\node[below] at (7.5,-0.2) {第4季};
\node[below] at (10,-0.2) {第2年};

% 活動標示
\node[above, align=center] at (1.25,0.5) {教師訓練\\EPA設計};
\node[above, align=center] at (3.75,0.5) {試行評估\\系統建置};
\node[above, align=center] at (6.25,0.5) {正式實施\\成效評估};
\node[above, align=center] at (8.75,0.5) {檢討改善\\擴大推廣};

% 里程碑
\foreach \x/\text in {2.5/完成設計, 5/開始實施, 7.5/初步成效, 10/全面推廣} {
  \fill[cardioBlue] (\x,0) circle (0.1);
  \node[rotate=45, anchor=west] at (\x,0.2) {\scriptsize \text};
}
\end{tikzpicture}
\end{center}
\end{frame}

\begin{frame}{成功實施的關鍵因素}
\begin{alertblock}{組織層面}
\begin{itemize}
\item 領導層支持與承諾
\item 充足資源投入
\item 跨部門協調合作
\end{itemize}
\end{alertblock}

\begin{block}{執行層面}
\begin{itemize}
\item 全體教師共識建立
\item 標準化訓練實施
\item 持續品質監控
\item 定期檢討改善
\end{itemize}
\end{block}

\begin{exampleblock}{文化層面}
\begin{itemize}
\item 學習型組織建立
\item 創新教育文化
\item 持續改善精神
\end{itemize}
\end{exampleblock}
\end{frame}

\section{討論與Q\&A}

\begin{frame}{討論重點}
\begin{enumerate}
\item EPA在內科各次專科的適用性如何？
\item 現有教學資源是否足夠支持EPA實施？
\item 評估者訓練的具體規劃為何？
\item 與現行評估制度如何整合？
\item 學生與住院醫師的接受度評估？
\end{enumerate}

\vspace{1cm}
\begin{center}
\Large{\textbf{歡迎提問與討論}}
\end{center}
\end{frame}

\begin{frame}{結語}
\begin{center}
\Large{\textcolor{cardioBlue}{\textbf{EPA不只是評估工具}}}
\end{center}

\vspace{0.5cm}

\begin{center}
\large{更是醫學教育轉型的契機}
\end{center}

\vspace{1cm}

\begin{itemize}
\item 從知識傳授到能力培養
\item 從標準化到個別化
\item 從單向教學到雙向學習
\item 從孤立評估到整合發展
\end{itemize}

\vspace{1cm}

% \begin{center}
% \textbf{讓我們攜手打造更優質的醫學教育！}
% \end{center}
% \end{frame}

\end{document}